\chapter{System Decomposition and Execution Model}

This chapter explains how the project maps the GlobalPlatform concepts
introduced in Chapter~4 onto an explicit software decomposition. In
other words, the previous chapter described \emph{what} GlobalPlatform
manages; this chapter describes \emph{how} the present system chooses to
organize that management internally.

%--------------------------------------------------
\section{Reading Chapter 4 as a Software Architecture}

% Plan:
% - explain that chapter 4 was conceptual, while chapter 5 is the chosen mapping
% - state explicitly that GP does not force a unique internal decomposition
% - introduce the guiding implementation idea

Chapter~4 distinguished three major conceptual roles:
\begin{itemize}
\item managed objects such as load files and installed applications;
\item administrative principals, namely Security Domains;
\item a command-dispatch context determined by selection and by the
  active communication state.
\end{itemize}

GlobalPlatform does not require one unique internal implementation of
those roles. The present project adopts a deliberately explicit reading:
the small set of truly platform-global commands is handled directly by a
kernel layer, while most other APDUs are dispatched to the currently
selected application. In that model, Security Domains are not reduced
to hard-coded special cases in every branch of the kernel. They are
modeled as applications resident on the secure element with elevated authority and access to
privileged services.

This choice is meant to preserve two useful properties:
\begin{itemize}
\item APDU dispatch stays uniform: after selection, commands are routed
  to an AID-associated entity through a common path;
\item authority remains explicit: only entities with the right
  privileges may invoke sensitive kernel services.
\end{itemize}

At the same time, this chapter must be read against Chapter~2. The
chosen decomposition is not judged only by elegance or modularity. It is
judged by whether it helps protect the assets and mitigate the threats
identified in the threat model:
\begin{itemize}
\item integrity of the application lifecycle;
\item confidentiality and integrity of cryptographic keys;
\item isolation between applications;
\item integrity of the internal registry;
\item control over administrative channels;
\item resistance to malicious internal applications;
\item non-bypassability of privileged operations.
\end{itemize}

The rest of the chapter therefore treats the software decomposition as a
security architecture, not merely as an implementation convenience.

%--------------------------------------------------
\section{Layered Decomposition}
\label{sec:system-layering}

% Plan:
% - present the security-relevant 3-layer view
% - mention core as supporting substrate
% - make clear that this is a software decomposition, not a GP normative statement

For the present project, the security-relevant execution stack is best
understood as a \textbf{three-layer architecture} resting on a reusable
runtime substrate.

At the bottom sits \textbf{core}, the reusable support layer. It
provides bootstrapping, target dispatch, serial transport, semihosting,
flash, cryptographic primitives, and similar low-level services. This
layer is essential to the implementation, but it is not by itself the
main authority boundary discussed in this chapter.

Above that substrate, the architecture of interest is:

\begin{itemize}
\item \textbf{kernel}: the platform core that owns secure-element-wide control
  flow, selection, privileged resource mediation, and the security
  boundary against applications;
\item \textbf{GP administrative applications}: first-class applications
  resident on the secure element, chiefly Security Domains, that implement administrative
  semantics but rely on privileged kernel services for sensitive effects;
\item \textbf{ordinary applications}: non-administrative application
  instances, typically loaded after issuance, which provide business
  logic but do not own management authority.
\end{itemize}

This decomposition should be read as an implementation architecture, not
as a claim that the GlobalPlatform specification itself mandates those
exact software layers. Its purpose is to realize the conceptual
distinctions of Chapter~4 while making the authority and isolation
boundaries analyzable from a security point of view.

The reason for making the three layers explicit is simple:
\begin{itemize}
\item the kernel must protect global assets and enforce non-bypassable
  mediation;
\item GP administrative applications must be powerful enough to express
  secure-element-management logic;
\item ordinary applications must remain strongly constrained so that a
  malicious ordinary application cannot compromise the rest of the
  platform.
\end{itemize}

%--------------------------------------------------
\section{Kernel Responsibilities}
\label{sec:system-kernel-responsibilities}

% Plan:
% - explain what stays in the kernel
% - emphasize OPEN-level handling and platform-global state
% - distinguish this from GP command semantics executed by security domains

In the present design, the kernel is intentionally small in semantic
scope. Its primary responsibilities are:

\begin{itemize}
\item boot and runtime entry;
\item transport-level APDU handling and selected-context maintenance;
\item direct handling of the small set of platform-global commands that
  must preempt ordinary dispatch;
\item enforcement of privileged operations that should not be performed
  by arbitrary applications;
\item ownership of platform-global resources such as registry storage,
  persistent management state, and protected cryptographic material.
\end{itemize}

At the APDU level, this means that the kernel directly handles the
\texttt{OPEN}-like part of the communication model. At the current stage
of the project, this is mainly \texttt{SELECT}; later it may include
\texttt{MANAGE CHANNEL} if logical channels are implemented. The
important point is that these commands affect the communication context
itself and therefore cannot simply be treated as ordinary APDUs routed to
the currently selected application.

Everything beyond that small preemptive surface is intended to follow the
normal selected-application dispatch path.

From a security perspective, the kernel is also where several Chapter~2
requirements become concrete:
\begin{itemize}
\item \textbf{non-bypassability}: all privileged effects must cross a
  kernel-controlled boundary;
\item \textbf{registry integrity}: authoritative updates to identity,
  privilege, lifecycle, and dependency state belong to kernel-protected
  code and data;
\item \textbf{control over administrative channels}: context-defining
  commands and secure-channel state must not be forgeable by ordinary
  applications;
\item \textbf{runtime integrity}: applications must not gain arbitrary
  code-execution power inside the kernel.
\end{itemize}

This is why the kernel is intentionally kept small in semantic scope. A
smaller kernel surface is easier to reason about for integrity and
non-bypassability than a monolithic design in which every management
command is implemented directly in the privileged core.

%--------------------------------------------------
\section{Administrative Applications and Security Domains}
\label{sec:system-admin-apps}

% Plan:
% - explain security domains as preinstalled privileged applications
% - connect to chapter 4 administration model
% - explain why GP commands can be handled there without making kernel huge

Chapter~4 explained that Security Domains are the administrative
principals of GlobalPlatform. In the present system decomposition, they
are realized as \textbf{administrative application-like objects}. They are
selected through AIDs like other resident entities on the secure element, carry
their own state, keys, privileges, and secure-channel context, and may be backed
either by kernel code or by an isolated Rustlet.

Oxide SE deliberately keeps the management hook on the kernel side. The kernel
recognizes the GlobalPlatform management commands that mutate platform state:
registry updates, lifecycle transitions, key installation, controlled content
loading, deletion, and secure-channel establishment. The active Security Domain
does not edit the registry directly. Instead, it authorizes or refuses the
operation, supplies the policy and cryptographic context, and may refine that
decision through its backend implementation.

This is the point where Oxide SE differs from a more Java Card-like
implementation strategy. A Java Card implementation may make the selected
Security Domain look like the functional APDU endpoint for many management
commands, with internal system services performing the protected effects behind
that endpoint. Oxide SE keeps the same visible administrative model, but uses a
symmetric internal decomposition: the kernel owns parsing, registry mutation,
non-escalation checks, and atomicity; the Security Domain owns authority,
policy, secure-channel material, and the decision to authorize management
families.

This point deserves to be stated strongly: in this design, a Security
Domain is not treated as a fake application or as a mere label attached to
kernel code. It is treated as a \textbf{first-class application-like object}.
It has an AID, may be selected, carries administrative state, and owns
privileges and key material associated with its role. In other words, the
project takes seriously the idea from Chapter~4 that a Security Domain is
simultaneously:
\begin{itemize}
\item an addressable entity resident on the secure element;
\item an administrative principal;
\item a holder of management authority and cryptographic trust anchors.
\end{itemize}

The security consequence is subtle but important. Because Security Domains may
be backed by isolated Rustlets in this architectural sense, they must also be
considered under the same isolation discipline as other application code. They
are more privileged, but they are not thereby exempt from needing a clear
boundary with the kernel and with peer applications.

This directly addresses two threats from Chapter~2:
\begin{itemize}
\item a \textbf{malicious internal application} must not be able to
  impersonate a Security Domain merely by reproducing an APDU surface;
\item a \textbf{compromised administrative application} must not obtain
  unlimited power beyond the privileges and kernel services explicitly
  granted to it.
\end{itemize}

The design aim is therefore not ``Security Domains are trusted, ordinary
applications are untrusted.'' The stronger and more useful statement is:
all application-level code is isolated; some selected entities are granted
administrative privileges, but those privileges remain bounded, mediated by the
kernel, and analyzable as Security-Domain authority rather than as raw access to
the registry.

%--------------------------------------------------
\section{Ordinary Applications}
\label{sec:system-ordinary-apps}

% Plan:
% - describe ordinary apps as selected entities without kernel privileges
% - distinguish application APDUs from GP administrative commands
% - prepare future loadable FAE model

Ordinary applications are the selected applications that implement business
logic rather than secure-element administration. Once selected, they receive APDUs
through the same general dispatch path as any other application-like
entity, but they do not own privileged management authority.

This distinction is what separates a true administrative Security Domain
from an ordinary application that merely imitates some of its APDU
surface. An ordinary application may expose commands that resemble
administrative ones, but without privileged kernel services it cannot
actually perform protected operations such as installing content,
changing lifecycle states, or replacing platform keys.

This is also the model that naturally extends to future loadable Rust
applications. They may be selected and invoked through APDUs, but their
rights remain subject to the authority structure imposed by the
administrative applications and enforced by the kernel.

This layer is where the \textbf{malicious internal application} threat of
Chapter~2 becomes central. An ordinary application must be assumed
capable of:
\begin{itemize}
\item sending malformed requests through any allowed inter-application or
  service interface;
\item attempting to read, corrupt, or confuse data belonging to other
  applications;
\item imitating management-level command formats in order to probe the
  authority boundary;
\item exploiting lifecycle, registry, or service-call mistakes to gain
  stronger rights than intended.
\end{itemize}

The architecture of this chapter is acceptable only if it helps contain
such a component. In particular, an ordinary application must not be able
to:
\begin{itemize}
\item alter registry state directly;
\item install or delete content directly;
\item manipulate platform key material directly;
\item forge the existence of a secure administrative session;
\item interfere with the selection state, memory, or persistent data of
  another application outside the explicitly allowed sharing mechanisms.
\end{itemize}

%--------------------------------------------------
\section{Selection and Dispatch Semantics}
\label{sec:system-selection-dispatch}

% Plan:
% - explain exact dispatch rule: SELECT preempts, rest goes to selected AID
% - make link with GP OPEN and selected-application semantics
% - explain GET RESPONSE stays transport-local and not a business dispatch decision

The dispatch rule adopted by the project is intentionally simple.

\begin{itemize}
\item \texttt{SELECT} preempts ordinary dispatch because it changes the
  selected context itself;
\item other future \texttt{OPEN}-level commands, such as
  \texttt{MANAGE CHANNEL}, belong to the same preemptive family if they
  are implemented;
\item every other business command is routed to the code associated with
  the currently selected application.
\end{itemize}

This means that a Security Domain and an ordinary application are, at the
dispatch level, treated in the same structural way: both are selected
entities addressed by APDUs after the context-setting phase. The
difference lies in what they are then allowed to do.

This is one of the most useful architectural simplifications of the
design. The routing rule remains uniform, but the authority rule does
not. The dispatch decision answers \emph{who receives the APDU}; the
security model separately answers \emph{what that receiver is allowed to
cause inside the platform}. Keeping these two questions distinct helps
avoid a common design pitfall in which routing shortcuts silently become
privilege shortcuts.

The transport exceptions described in Chapter~8 remain orthogonal to this
business dispatch rule. For instance, \texttt{GET RESPONSE} is handled by
the transport loop because it belongs to the \texttt{T=0} realization of
the exchange rather than to the business semantics of the selected
application.

From the point of view of Chapter~2, this section is where the security
goal of \textbf{control over administrative channels} meets the execution
model. The selected context must be explicit, coherent, and protected
against spoofing or unintended reuse. Otherwise, an attacker could move
from malformed APDU traffic to actual misbinding of management authority.

%--------------------------------------------------
\section{Privileged Services and Authority Boundaries}
\label{sec:system-authority-boundaries}

% Plan:
% - explain why kernel services must mediate sensitive actions
% - explain capability idea without over-specifying syscall ABI
% - tie back to chapter 4 privileges and administrative model

The authority model of Chapter~4 reappears here as a kernel-service
boundary. Sensitive state changes are not performed merely because an
application emitted an APDU handler branch claiming to be administrative.
They are performed only when the selected application is recognized as holding
the corresponding authority and when it invokes the privileged kernel
operation that implements the requested effect.

This allows a clean separation between:

\begin{itemize}
\item \textbf{APDU semantics}, which may be expressed in the code of an
  administrative application;
\item \textbf{effective authority}, which remains enforced by the kernel.
\end{itemize}

The future syscall or service boundary of the project should therefore be
understood as the concrete mechanism by which Security Domains ask the
kernel to perform privileged GlobalPlatform actions. Ordinary
applications may use unprivileged services, but they should not be able
to cross that authority boundary by merely reproducing the visible APDU
format of a management command.

This service boundary is also the place where the architecture should
remain most explicit about privilege granularity. A Security Domain
should not receive a generic ``kernel mode'' escape hatch. It should
receive only the specific privileged operations needed to realize its
GlobalPlatform role. This is important for two reasons:
\begin{itemize}
\item it limits the blast radius of a compromised administrative
  application;
\item it keeps the privilege model aligned with the GlobalPlatform object
  model of Chapter~4 instead of collapsing everything into one
  undifferentiated superuser privilege.
\end{itemize}

In that sense, the service boundary is not only an implementation device;
it is one of the principal security controls of the system.

%--------------------------------------------------
\section{Isolation Model}
\label{sec:system-isolation-model}

% Plan:
% - make explicit the different isolation levels
% - connect them to threat model
% - distinguish kernel/application, admin/ordinary, and app/app isolation

The software stack just described naturally induces several nested
isolation relations. Making them explicit is important, because
Chapter~2 identified \textbf{isolation between applications},
\textbf{runtime integrity}, and \textbf{non-bypassability} as core
security goals.

\subsection{Kernel-versus-Application Isolation}

The strongest boundary is the one separating the kernel from every
application-level component, whether administrative or ordinary. The
kernel owns the trusted execution substrate for:
\begin{itemize}
\item selected-context management;
\item registry and lifecycle authority;
\item key ownership and secure-channel state;
\item privileged effects on persistent platform state.
\end{itemize}

Applications may request services across this boundary, but they must not
be able to bypass it. In particular, neither a Security Domain nor an
ordinary application should be able to transform ``being selected'' into
``executing with arbitrary kernel authority.'' This boundary is the main
answer to the Chapter~2 assumptions about trusted execution and runtime
integrity.

\subsection{Administrative-Application Isolation}

Administrative applications, notably Security Domains, sit above the
kernel but below ordinary applications in authority. They are privileged
with respect to the service interface, but they are still applications.
They therefore require their own containment properties:
\begin{itemize}
\item they should not be able to tamper with kernel memory or code;
\item they should not automatically inherit each other's authority if
  multiple Security Domains coexist;
\item they should only access platform state through the kernel-mediated
  interfaces and the privileges explicitly granted to them.
\end{itemize}

This is the right way to think about a compromised administrator in the
present architecture. Compromise of an administrative application is
serious, but it should still be analyzable in terms of the privileges
granted to that Security Domain and the kernel operations it can
legitimately invoke.

\subsection{Ordinary-Application Isolation}

Ordinary applications require mutual isolation from one another. At a
minimum, this means that one application should not be able to:
\begin{itemize}
\item read or overwrite another application's private state;
\item alter another application's lifecycle or registration state;
\item hijack another application's selected context;
\item exploit a shared service path to escape its own authority level.
\end{itemize}

This is the direct architectural response to the malicious internal
application adversary of Chapter~2. The fact that applications share the
same APDU transport and may coexist under the same Security Domain must
not be allowed to collapse their isolation boundaries.

\subsection{Administrative-versus-Ordinary Isolation}

One particularly important case is the isolation boundary between
administrative applications and ordinary ones. The former may legitimately
manage the latter, but that management relationship must not be confused
with arbitrary data or code-level visibility.

Put differently: a Security Domain may hold authority \emph{over} an
application in the sense of installation, lifecycle, privilege
assignment, personalization, or deletion, without thereby becoming a
mechanism that breaks all runtime isolation with that application. This
distinction is important if the platform later wants to support several
administrative domains or to reason about delegated management without
collapsing all application boundaries.

\subsection{Isolation Summary}

The intended hierarchy can therefore be summarized as follows:
\begin{itemize}
\item the \textbf{kernel} is isolated from all applications and mediates
  all privileged effects;
\item \textbf{GP administrative applications} are isolated from the
  kernel, but possess privileged, kernel-mediated capabilities;
\item \textbf{ordinary applications} are isolated from the
  kernel and from GP administrative applications except through explicit
  mediated interfaces;
\item \textbf{applications in general} are isolated from one another,
  whether they are ordinary or administrative, except where the platform
  intentionally defines a controlled sharing or delegation mechanism.
\end{itemize}

This is the main security reading of the three-layer architecture. It is
what allows the chapter to claim that Security Domains can be treated as
full applications without giving up the isolation goals that motivated
the architecture in the first place.

%--------------------------------------------------
\section{Discussion}

% Plan:
% - summarize chapter 5 as the system reading of chapter 4
% - point forward to implementation chapter

This chapter has translated the conceptual distinctions of Chapter~4 into
a software architecture: a small kernel handles context-defining
commands, selected applications receive ordinary APDUs through a common
dispatch mechanism, Security Domains are realized as privileged
administrative applications, and sensitive effects remain mediated by the
kernel.

This is the architectural reading that guides the reference
implementation. Chapter~8 returns to it from a code-level perspective and
describes how the transport loop, dispatch layer, cryptographic support,
and future registry and lifecycle services realize this decomposition in
practice.

Its main security claim is not that privileged applications disappear,
but that they can be treated as first-class applications without
sacrificing analyzable authority boundaries. The kernel remains the
non-bypassable mediation point; Security Domains remain powerful but
bounded administrative principals; and ordinary applications remain
isolated both from the kernel and from one another.

%--------------------------------------------------
